% Title page for the AJSE submission (uploaded separately as "Title Page").
% AJSE: title, author information (affiliation, corresponding author e-mail, ORCID), abstract,
% keywords, acknowledgements and Statements and Declarations belong on the title page.
\documentclass[11pt,a4paper]{article}
\usepackage[a4paper,margin=2.2cm]{geometry}
\usepackage[T1]{fontenc}
\usepackage[utf8]{inputenc}
\usepackage{lmodern}
\usepackage{microtype}
\usepackage{xcolor}
\usepackage[hidelinks]{hyperref}
\newcommand{\ablPrimBurDMMax}{0.84}
\newcommand{\ablPrimBurDMMin}{0.65}
\newcommand{\ablPrimBurLRMax}{0.63}
\newcommand{\ablPrimBurLRMin}{0.58}
\newcommand{\ablPrimBurOSMax}{0.69}
\newcommand{\ablPrimBurOSMin}{0.65}
\newcommand{\ablPrimVolDMMax}{0.78}
\newcommand{\ablPrimVolDMMin}{0.65}
\newcommand{\ablPrimVolLRMax}{0.63}
\newcommand{\ablPrimVolLRMin}{0.58}
\newcommand{\ablPrimVolOSMax}{0.70}
\newcommand{\ablPrimVolOSMin}{0.65}
\newcommand{\ablTotBurDMMax}{0.87}
\newcommand{\ablTotBurDMMin}{0.85}
\newcommand{\ablTotBurLRMax}{0.63}
\newcommand{\ablTotBurLRMin}{0.59}
\newcommand{\ablTotBurOSMax}{0.73}
\newcommand{\ablTotBurOSMin}{0.70}
\newcommand{\ablTotVolDMMax}{0.88}
\newcommand{\ablTotVolDMMin}{0.85}
\newcommand{\ablTotVolLRMax}{0.63}
\newcommand{\ablTotVolLRMin}{0.58}
\newcommand{\ablTotVolOSMax}{0.73}
\newcommand{\ablTotVolOSMin}{0.72}
\newcommand{\clinOnlyDMMax}{0.81}
\newcommand{\clinOnlyDMMin}{0.79}
\newcommand{\clinOnlyLRMax}{0.57}
\newcommand{\clinOnlyLRMin}{0.51}
\newcommand{\clinOnlyOSMax}{0.70}
\newcommand{\clinOnlyOSMin}{0.69}
\newcommand{\clsAvail}{1}
\newcommand{\clsClinDMMax}{0.83}
\newcommand{\clsClinDMMin}{0.77}
\newcommand{\clsClinLRMax}{0.76}
\newcommand{\clsClinLRMin}{0.71}
\newcommand{\clsClinOSMax}{0.67}
\newcommand{\clsClinOSMin}{0.64}
\newcommand{\clsDMCVClin}{0.86}
\newcommand{\clsDMCVImg}{0.87}
\newcommand{\clsDMCohClin}{0.88}
\newcommand{\clsDMCohImg}{0.82}
\newcommand{\clsDMMax}{0.88}
\newcommand{\clsDMMin}{0.82}
\newcommand{\clsLRCVClin}{0.68}
\newcommand{\clsLRCVImg}{0.58}
\newcommand{\clsLRCohClin}{0.75}
\newcommand{\clsLRCohImg}{0.54}
\newcommand{\clsLRMax}{0.75}
\newcommand{\clsLRMin}{0.54}
\newcommand{\clsOSCVClin}{0.74}
\newcommand{\clsOSCVImg}{0.73}
\newcommand{\clsOSCohClin}{0.73}
\newcommand{\clsOSCohImg}{0.73}
\newcommand{\clsOSMax}{0.74}
\newcommand{\clsOSMin}{0.73}
\newcommand{\cnnDMCVClin}{0.88}
\newcommand{\cnnDMCVImg}{0.83}
\newcommand{\cnnDMCohClin}{0.93}
\newcommand{\cnnDMCohImg}{0.89}
\newcommand{\cnnLRCVClin}{0.57}
\newcommand{\cnnLRCVImg}{0.53}
\newcommand{\cnnLRCohClin}{0.59}
\newcommand{\cnnLRCohImg}{0.45}
\newcommand{\cnnOSCVClin}{0.68}
\newcommand{\cnnOSCVImg}{0.63}
\newcommand{\cnnOSCohClin}{0.69}
\newcommand{\cnnOSCohImg}{0.67}
\newcommand{\eTestDM}{7}
\newcommand{\eTestLR}{30}
\newcommand{\eTestOS}{45}
\newcommand{\figGraphK}{20}
\newcommand{\mainCIDMCohClin}{0.68--0.98}
\newcommand{\mainCIDMCohImg}{0.71--0.99}
\newcommand{\mainCILRCohClin}{0.46--0.72}
\newcommand{\mainCILRCohImg}{0.51--0.74}
\newcommand{\mainCIOSCohClin}{0.58--0.80}
\newcommand{\mainCIOSCohImg}{0.61--0.82}
\newcommand{\mainDMCVClin}{0.87}
\newcommand{\mainDMCVImg}{0.85}
\newcommand{\mainDMCohClin}{0.85}
\newcommand{\mainDMCohImg}{0.87}
\newcommand{\mainDMMax}{0.87}
\newcommand{\mainDMMin}{0.85}
\newcommand{\mainLRCVClin}{0.59}
\newcommand{\mainLRCVImg}{0.62}
\newcommand{\mainLRCohClin}{0.60}
\newcommand{\mainLRCohImg}{0.63}
\newcommand{\mainLRMax}{0.63}
\newcommand{\mainLRMin}{0.59}
\newcommand{\mainOSCVClin}{0.72}
\newcommand{\mainOSCVImg}{0.73}
\newcommand{\mainOSCohClin}{0.70}
\newcommand{\mainOSCohImg}{0.72}
\newcommand{\mainOSMax}{0.73}
\newcommand{\mainOSMin}{0.70}
\newcommand{\matchedN}{1944}
\newcommand{\matchedTestDeltaDM}{$-$0.006}
\newcommand{\matchedTestDeltaLR}{$-$0.002}
\newcommand{\matchedTestDeltaOS}{$-$0.010}
\newcommand{\matchedTestFracAll}{29}
\newcommand{\matchedTestFracDM}{26}
\newcommand{\matchedTestFracLR}{44}
\newcommand{\matchedTestFracOS}{14}
\newcommand{\matchedTestMedAll}{$-$0.006}
\newcommand{\matchedValFracAll}{34}
\newcommand{\matchedValFracDM}{22}
\newcommand{\matchedValFracLR}{55}
\newcommand{\matchedValFracOS}{34}
\newcommand{\mlpDMCVClin}{0.86}
\newcommand{\mlpDMCVImg}{0.85}
\newcommand{\mlpDMCohClin}{0.86}
\newcommand{\mlpDMCohImg}{0.86}
\newcommand{\mlpLRCVClin}{0.59}
\newcommand{\mlpLRCVImg}{0.59}
\newcommand{\mlpLRCohClin}{0.58}
\newcommand{\mlpLRCohImg}{0.64}
\newcommand{\mlpOSCVClin}{0.72}
\newcommand{\mlpOSCVImg}{0.74}
\newcommand{\mlpOSCohClin}{0.72}
\newcommand{\mlpOSCohImg}{0.73}
\newcommand{\nPatCan}{294}
\newcommand{\nPatMAA}{137}
\newcommand{\nTestDM}{109}
\newcommand{\nTestLR}{121}
\newcommand{\nTestOS}{135}
\newcommand{\nodalFracMedCan}{44}
\newcommand{\nodalFracMedMAA}{30}
\newcommand{\nodalPctCan}{37}
\newcommand{\nodalPctMAA}{59}
\newcommand{\pairedDeltaMax}{+0.050}
\newcommand{\pairedDeltaMin}{$-$0.023}
\newcommand{\pairedMinP}{0.03}
\newcommand{\pairedN}{18}
\newcommand{\pairedNsig}{1}
\newcommand{\seedRepeats}{10}
\newcommand{\seedSDmax}{0.004}
\newcommand{\seedSDmed}{0.001}
\newcommand{\seedStableLoss}{3}
\newcommand{\seedStableWins}{9}
\newcommand{\winsCls}{9}
\newcommand{\winsGraph}{9}
\newcommand{\winsMLP}{8}
\newcommand{\winsPrimBur}{7}
\newcommand{\winsPrimVol}{6}
\newcommand{\winsTotBur}{9}
\newcommand{\winsTotVol}{9}

\setlength{\parindent}{0pt}
\setlength{\parskip}{6pt}
\begin{document}

\begin{center}
{\large\bfseries Graph neural networks on primary and nodal tumour radiomics for outcome prediction in head and neck cancer: a multicentre external validation study}\\[10pt]
{\normalsize Hasan Shaikh$^{*}$}\\[4pt]
{\small Department of Computer Science and Engineering, Aligarh Muslim University, Aligarh, Uttar Pradesh 202002, India}\\[6pt]
{\small $^{*}$Corresponding author. E-mail: \href{mailto:hasanshaikh3198@gmail.com}{hasanshaikh3198@gmail.com}}
% ORCID is optional for AJSE ("if available"). To add it, insert a line such as:
% \\[2pt]{\small ORCID: 0000-0000-0000-0000}
\end{center}

\section*{Abstract}
Graph neural networks that connect each patient to similar patients are increasingly used to predict outcomes from computed tomography (CT) radiomics, yet the graph's own contribution is seldom tested. This study compared such networks with a published convolutional neural network (CNN) for head and neck squamous cell carcinoma on its public benchmark protocol. Pretreatment CT scans of \nPatCan{} patients from four Canadian centres were used for development and \nPatMAA{} patients from a Dutch centre for external testing. Endpoints were locoregional recurrence and distant metastasis at two years and death at four years. The proposed network describes the primary and nodal tumour together and links each patient to similar patients. Graph convolutional, GraphSAGE and graph attention versions were each compared with an identical network without edges, and a factorial ablation separated tumour definition from nodal-burden clinical variables. This whole-tumour network reached external areas under the receiver operating characteristic curve of \mainLRMin--\mainLRMax{} for recurrence, \mainDMMin--\mainDMMax{} for metastasis and \mainOSMin--\mainOSMax{} for survival. It outperformed the CNN in \winsGraph{} of 12 settings, compared with \winsPrimVol{} of 12 for a primary-tumour variant, and including the nodes raised the metastasis value by up to 0.21. Removing the edges changed little: the edgeless network outperformed the CNN in \winsMLP{} of 12 settings, and only \pairedNsig{} of \pairedN{} paired differences was significant, in its favour. The gain came from describing the whole tumour, and graph-based outcome models should be reported alongside their edgeless counterpart.

\textbf{Keywords:} Head and neck cancer, Radiomics, Graph neural network, Nodal tumour volume, External validation, Radiotherapy outcome prediction

\section*{Acknowledgements}
The author thanks the investigators who made the Head-Neck-PET-CT and HN1 collections publicly available through The Cancer Imaging Archive.

\section*{Statements and Declarations}

\paragraph{Funding}
No funding was received for conducting this study.

\paragraph{Competing Interests}
The author has no relevant financial or non-financial interests to disclose.

\paragraph{Author Contributions}
Conceptualization, methodology, software, formal analysis, investigation, data curation, visualization, writing (original draft) and writing (review and editing): Hasan Shaikh. The author read and approved the final manuscript.

\paragraph{Data Availability}
All data analysed in this study are publicly accessible. The CT images, radiotherapy structure sets and clinical data of the development cohort are available from The Cancer Imaging Archive as the Head-Neck-PET-CT collection (\url{https://doi.org/10.7937/K9/TCIA.2017.8oje5q00}) and those of the external test cohort as the HN1 collection (\url{https://doi.org/10.7937/tcia.2019.8kap372n}), under the data usage policies of The Cancer Imaging Archive. No new patient data were generated. The radiomic features, tumour volumes and predictions derived in this study can be regenerated from these collections with the released code. The performance of the CNN was taken from the published report of Mateus et al. (Mateus et al., Sci. Rep. 2023).

\paragraph{Code Availability}
The code for feature extraction, construction of the patient-similarity graphs, model training and selection, the ablation and seed analyses, and the script that generates every table, figure and quoted number of this manuscript is openly available at \url{https://github.com/hash123shaikh/Radiomics_Graph_Attention_Network_HNC} (directory \texttt{grafs-gat/popgraph}) under the MIT licence, together with the software environment and step-by-step instructions to reproduce the results.

\section*{Compliance with Ethical Standards}

\paragraph{Disclosure of Potential Conflicts of Interest}
The author declares that there are no conflicts of interest.

\paragraph{Research Involving Human Participants and/or Animals}
This study is a secondary analysis of de-identified, publicly available data from The Cancer Imaging Archive; no new data were collected from human participants and no animals were involved. The original studies that created the collections obtained institutional ethics approval, and separate ethics approval was therefore not required for this analysis.

\paragraph{Informed Consent}
Not applicable to this secondary analysis. Informed consent, or a waiver of consent, was obtained by the original studies under their institutional approvals.

\end{document}
